\scanpage{153}
\mspage{69}
% archive sheet number: 152

\uline{Proposition} 4.3.4. Soit $M$ une $\unsure{U|\mathcal U}$-catégorie. Les conditions suivantes sont équivalentes

a) $M$ est une catégorie $U$-pseudo-topos, i.e. $M$ est stable par $U$-limites inductives, et admet une \inserted{$U$-}petite ss-catégorie \inserted{$S$ str.} génératrice.

b) Il existe une \deleted{\unsure{famille|foncteur}} $U$-petite catégorie $S$, et un foncteur $\unsure{i}:M\hookrightarrow S^\wedge$ \uline{pleinement fidèle}, ayant un adjoint à gauche.

\uline{Dém} On a $a\Rightarrow b$, car il suffit de prendre pour $S$ \inserted{dans b)} celui spécifié dans a), l'hyp. sur $S$ implique que le foncteur canonique $M\unsure{\longrightarrow|-}\deleted{\unsure{M}}S^\wedge$ composé de $M\mathrel{\mathop{\hookrightarrow}\limits^{\deleted{\unsure{h}}}}M^\wedge\mathrel{\mathop{\unsure{-|\longrightarrow}}\limits^{\mathrm{restr.}}}S^\wedge$ est pleinement fidèle. Qu'il admette un adj. à g. résulte du lemme 4.3.2.

L'implication $b)\Rightarrow a)$ résulte du lemme 4.3.3.

\uline{Corollaire} \inserted{4.3.5.} Soit $M$ une $U$-catégorie. Pour que $M$ soit un $U$-pseudo-topos \deleted{\unsure{$\zeta|\xi$}} \uline{accessible} (i.e \unsure{tous|tout} \unsure{dont|dans} les objets sont accessibles), il f. et s. qu'il existe $S$ dans $\mathrm{Cat}$ et un foncteur pleinement fidèle et \uline{accessible} $M\mathrel{\mathop{\unsure{\lhook\joinrel-|\hookrightarrow}}\limits^\alpha}S^\wedge$, admettant un adjoint à gauche $\rho$.

